{
\small
\renewcommand{\arraystretch}{1.15}
\setlist[enumerate]{leftmargin=*, nosep, labelsep=0.3em}

\begin{tabularx}{\textwidth}{|>{\raggedright\arraybackslash}p{3.5cm}|>{\raggedright\arraybackslash}X|}
\hline
\textbf{Tên Use case} & Tạo bài thi \\
\hline
\textbf{ID} & UC-0015 \\
\hline
\textbf{Tác nhân} & Người kiểm duyệt (Mod) hoặc Quản trị viên (Admin) \\
\hline
\textbf{Mô tả} & Người kiểm duyệt hoặc quản trị viên tạo bài thi mới trong hệ thống \\
\hline
\textbf{Điều kiện tiên quyết} & Người dùng đã đăng nhập và có vai trò Mod hoặc Admin \\
\hline
\textbf{Điều kiện sau} & Bài thi mới được tạo ở trạng thái nháp \\
\hline
\textbf{Kích hoạt} & Người dùng chọn Tạo bài thi mới trong trang quản lý \\
\hline
\textbf{Luồng chính} &
\begin{enumerate}
    \item Người dùng mở trang quản lý bài thi.
    \item Người dùng nhấn Tạo bài thi mới.
    \item Người dùng nhập tiêu đề, mô tả và thời gian làm bài.
    \item Hệ thống tạo bài thi ở trạng thái nháp và trả về ID bài thi.
    \item Người dùng tiếp tục thêm các phần (UC-0019), câu hỏi (UC-0020), và thẻ phân loại.
\end{enumerate} \\
\hline
\textbf{Luồng thay thế} &
\begin{enumerate}
    \item Không có luồng thay thế.
\end{enumerate} \\
\hline
\textbf{Luồng ngoại lệ} &
\begin{enumerate}
    \item \textbf{Không có quyền:}
    \begin{enumerate}
        \item[2Err.] Hệ thống trả về lỗi 403 khi người dùng không có vai trò Mod hoặc Admin.
    \end{enumerate}
    \item \textbf{Dữ liệu không hợp lệ:}
    \begin{enumerate}
        \item[4Err.] Hệ thống trả về lỗi khi thiếu các trường bắt buộc (tiêu đề).
    \end{enumerate}
\end{enumerate} \\
\hline
\end{tabularx}
\captionof{table}{Đặc tả Use case: UC-0015 Tạo bài thi}
}
